%% ── RELATED WORK ────────────────────────────────────────────────────────────
% Target: ≤0.5 page CoRL format (~300 words)

\section{Related Work}
\label{sec:related}

\paragraph{Vision-Language-Action models.}
Recent VLA models have achieved remarkable in-distribution performance on
tabletop manipulation benchmarks: \mbox{ACoT-VLA}~\cite{acot_vla_2026}
reaches 98.5\% on LIBERO via action chain-of-thought reasoning;
StarVLA~\cite{starvla_2026} provides a modular framework across four
action heads (FAST, OFT, flow-matching, DIT) achieving 97.8\%;
OpenVLA-OFT~\cite{openvla_oft_2025} improves to 97.1\% via parallel
decoding and action chunking.
Despite these gains, all models degrade \textbf{14–18\%} under spatial
distribution shifts~\cite{openvla_oft_2025,starvla_2026,acot_vla_2026}:
OpenVLA-OFT falls to 79.6\%, StarVLA to 79.7\%, ACoT-VLA to 84.1\%.
This degradation is not a capacity limitation — it is \emph{frame dependence}:
actions represented in absolute world or camera coordinates change completely
under any rigid-body transform of the workspace.
VLA-0~\cite{vla0_2025} confirms that even minimal text-based action
discretization achieves 94.7\% in-distribution, yet inherits the same
brittleness to spatial shifts.

\paragraph{Action tokenization.}
\fast{}~\cite{fast2025} achieves $\sim$10$\times$ compression via DCT + BPE
and enables dexterous tasks where per-dimension binning fails.
\oat{}~\cite{oat2026} introduces ordered token spaces via register encoders,
finite scalar quantization, and nested dropout, enabling prefix-decodable
autoregressive policies.
Both operate on raw action trajectories in their original coordinate frames.
A trajectory translated or rotated in space produces entirely different
tokens — neither method addresses geometric invariance.

\paragraph{Flow-matching policies.}
$\pi_0$/$\pi_{0.5}$~\cite{pi0_2024} achieve impressive dexterity via
continuous flow matching at 50\,Hz, avoiding discretization.
The probabilistic nature of diffusion provides some robustness, but does
not produce SE(3)-invariant action representations.

\paragraph{3D-aware VLA and data augmentation.}
3D-VLA and SpatialVLA encode 3D structure in visual representations, yet
still regress actions in world-frame coordinates --- geometrically-aware
perception does not imply invariant actions.
A common practical solution is SE(3) data augmentation: apply random rigid
transforms to training demonstrations. While this improves in-distribution
diversity, the \emph{token stream itself} remains frame-dependent (BPE tokens
change 99.1\% under SE(3), even with augmented training data), and
out-of-distribution spatial configurations still require re-collection.
\dhb{} provides invariance \emph{by construction}, not by augmentation.

\paragraph{SE(3)-equivariant manipulation policies.}
EquivAct~\cite{equivact2024} and RISE~\cite{rise2024} achieve SE(3)-equivariant
policies via equivariant network architectures, demonstrating strong
real-robot spatial generalization.
These works build equivariance into the \emph{policy architecture};
\dhb{} builds invariance into the \emph{action representation}.
The approaches are orthogonal: \dhb{}'s invariant tokens could serve as
action space for equivariant policies.

\paragraph{Spatial structure and invariant representations.}
3D-VLA and SpatialVLA encode 3D structure in visual representations, yet
still regress actions in world-frame coordinates.
The DHB framework~\cite{soloperto2015,lee2018} established bidirectional
invariant trajectory encoding on SE(3); Rotation-Minimizing Frames~\cite{wang2008rmf}
via Householder reflections provide smooth frame propagation.
\dhb{} extends these foundations — which are not themselves novel — to VLA
tokenization: our contribution is the combination of invariant encoding
with an OAT-style tokenizer (register encoder + FSQ + nested dropout,
reimplemented from Liu et al.~\cite{oat2026}) applied to invariant rather
than raw action sequences, enabling geometric adaptation at test time.
